\chapter{Geltungsbereich und Paarbegriffe}
\label{b14guide:ch:foundations}

Dieser Leitfaden ergänzt Band~14 um eine ausgewählte Übersicht der
wichtigsten Definitionen, Charakterisierungen und unmittelbaren
Folgerungen zu \emph{Paarinfima und Paarsuprema}. Er enthält keine
Beweise oder Beweisskizzen und erhebt keinen Anspruch auf eine
vollständige Liste aller bekannten oder ableitbaren Sätze.
Die Nummern mit Präfix $14H$ gehören ausschließlich zu dieser Ergänzung.

\section{Voraussetzungen und Bandgrenzen}

Durchgehend ist $(P,\leq)$ eine halbgeordnete Menge im Sinne von
Band~12. Alle Elementvariablen liegen in $P$, sofern kein anderer
Träger angegeben ist. $P$ darf leer sein. Die mengentheoretische
Grundlage ist ZF ohne zusätzliches Auswahlaxiom.

Gegenstand sind ausschließlich die Werte zu einzelnen Paaren
$\{x,y\}$ und die daraus gebildeten Ausdrücke.
Es wird \emph{nicht} vorausgesetzt, daß jedes Paar ein Infimum
oder Supremum besitzt. Auch $x=y$ ist zugelassen.
Die Existenzbedingungen in einem Satz betreffen nur die dort
genannten Paare; sie sind keine zusätzlichen Strukturaxiome.

\begin{center}
\small
\begin{tabularx}{\linewidth}{@{}p{.72\linewidth}X@{}}
\toprule
Eigenständiges Thema außerhalb dieses Leitfadens & Band\\
\midrule
Allgemeine Funktionen; Injektivität, Surjektivität, Bijektivität & 5--8\\
Äquivalenzrelationen & 11\\
Allgemeine Theorie der Halbordnungen & 12\\
Schranken, Infima und Suprema beliebiger Teilmengen & 13\\
Totale Ordnungen und Wohlordnungen & 15--16\\
Halbverbände und Verbände als axiomatische Strukturen & 45\\
\bottomrule
\end{tabularx}
\end{center}

Die Paaroperationen werden hier durch ihre Ordnungseigenschaft
beschrieben. Allgemeine Funktionensätze, Homomorphismen,
Darstellungs- und Vollständigkeitssätze sowie Theorien mit
Distributivitäts-, Modularitäts- oder Komplementaxiomen werden
nicht entwickelt. Die Mengenfamilien in Kapitel~\ref{b14guide:ch:sets}
realisieren die Paarbegriffe des Hauptbands.

\section{Definition und Existenz}

\begin{definition}[Paarinfimum und Paarsupremum]
\label{b14guide:def:pair}
Für $x,y\in P$ heißt $m\in P$ \emph{Paarinfimum von $x,y$}, wenn
\[
 m\leq x,\quad m\leq y,\quad
 \forall z\in P\;\bigl((z\leq x\land z\leq y)\Rightarrow z\leq m\bigr).
\]
Ein $s\in P$ heißt \emph{Paarsupremum von $x,y$}, wenn
\[
 x\leq s,\quad y\leq s,\quad
 \forall z\in P\;\bigl((x\leq z\land y\leq z)\Rightarrow s\leq z\bigr).
\]
Es handelt sich um das Infimum beziehungsweise Supremum der
ungeordneten Paarmenge $\{x,y\}$ im vorgegebenen Träger $P$.
\end{definition}

\begin{theorem}[Eindeutigkeit der Paarwerte]
\label{b14guide:thm:unique}
Für jedes Paar $x,y\in P$ gibt es höchstens ein Paarinfimum und
höchstens ein Paarsupremum.
\end{theorem}

\begin{definition}[Existenzprädikate und Schreibweise]
\label{b14guide:def:existence}
$\InfEx(x,y)$ bedeutet, daß das Paarinfimum von $x,y$ in $P$
existiert; $\SupEx(x,y)$ bedeutet die entsprechende Existenz
des Paarsupremums. Im jeweiligen Existenzfall schreiben wir
\[
 x\wedge y=\inf_{P,\leq}\{x,y\},\qquad
 x\vee y=\sup_{P,\leq}\{x,y\}.
\]
Die Existenzprädikate entsprechen
$\operatorname{InfEx}_{P,\leq}$ und $\operatorname{SupEx}_{P,\leq}$
des Hauptbands. Bei mehreren Trägern kennzeichnen
$\wedge_P,\vee_P$ den jeweiligen Bezug.
\end{definition}

Eine Gleichung wie $x\wedge y=m$ schließt die Existenz des
linken Wertes ein. Bei Ungleichungen und verschachtelten
Ausdrücken nennen die Sätze die erforderliche Existenz ausdrücklich.
Es gibt keine Ersatzwerte für nicht existierende Paarwerte.
In $\InfEx(x\vee y,z)$ muß insbesondere $x\vee y$ bereits
vorhanden sein. Die Zeichen $\wedge,\vee$ stehen zwischen
Elementen für Paaroperationen, zwischen Aussagen für
logische Konjunktion und Disjunktion.